\section{Mapping-cone squares and a module extension}\label{sec:cones}

The mapping cone of an actual sphere map supplies the realization
information needed by the algebraic obstruction. We first recover a
cone square from the square Hurewicz image of the loop adjoint.
For targets \(QS^n\) with \(n>0\), we then retain the top cone class
while passing to indecomposable cohomology, and obtain a nonzero tensor
in \(I\otimes_A M_k(n)\). The first step is valid for every connected
target, including those that are not simply connected.

\begin{lemma}[The cone-square detector]\label{lem:cone-square}
Let \(T\) be a connected based space and \(m\ge2\). Suppose that
\(g:S^{2m-1}\to T\) has zero mod-2 Hurewicz image and that its adjoint
\(g^\flat:S^{2m-2}\to\Omega T\) has Hurewicz image \(x^2\), where
\(x\in H_{m-1}(\Omega_0T;\F_2)\) and the square uses the loop product.
If \(u\in H^m(T;\F_2)\) satisfies \(u^2=0\), then it has a unique
extension \(\widetilde u\in H^m(C_g;\F_2)\), and
\begin{equation}\label{eq:cone-square}
\widetilde u^2=\langle u,\sigma_*x\rangle a,
\end{equation}
where \(a\in H^{2m}(C_g;\F_2)\) is the image of the top relative class.
This relative class is nonzero.
\end{lemma}

\begin{proof}
The cohomology sequence of the attachment shows that restriction in
degree \(m\) is an isomorphism: \(m<2m-1\). It also shows that the
top relative class injects, since \(g^*:H^{2m-1}(T)\to H^{2m-1}(S^{2m-1})\)
is zero. Thus \(\widetilde u^2\) is a scalar multiple of \(a\).

To determine the scalar, we use the universal space for a cohomology
class equipped with a nullhomotopy of its square. Write
\(K(j)=K(\F_2,j)\) and set
\[
F=\hofib\bigl(\Sq^m:K(m)\longrightarrow K(2m)\bigr).
\]
A map representing \(u\) lifts to \(F\) because its square is zero.
Looping the defining map represents \(\Sq^m\) on the fundamental class
of degree \(m-1\), which is zero by instability. Consequently the
fibration
\[
K(2m-2)\longrightarrow\Omega F\longrightarrow K(m-1)
\]
is trivial up to fibre homotopy. This is a statement about spaces over
\(K(m-1)\); its splitting need not preserve loop products.

We need the additive square-zero property of Eilenberg--MacLane spaces:
every positive-degree mod-two homology class of \(K(j)\), \(j\ge1\),
has zero square for its additive Pontryagin product. This follows from
Serre's cohomology calculation \cite{Serre1953}; the following chain
argument gives the precise property used here. Use the simplicial
abelian group \(G_r=Z^j(\Delta[r];\F_2)\) of normalized simplicial
\(j\)-cocycles, with faces given by restriction. Its Moore complex is
\(\F_2\) in degree \(j\) and zero elsewhere: below that degree there
are no \(j\)-cochains; in higher degrees the vanishing faces and the
cocycle equation force the remaining face to vanish. Thus its realization
is an additive model for \(K(j)\). The normalized \emph{linearized}
complex \(N(\F_2[G])\), rather than the Moore complex of \(G\),
computes its homology. Addition induces the chain product through the
shuffle map and \([a][b]=[a+b]\). For a chain \(z\) of degree \(q>0\),
the terms in \(z*z\) pair under exchange of the two complementary
\((q,q)\)-shuffle blocks and of the two chosen simplices of \(z\).
This involution has no fixed point, and the paired summands agree because
addition is commutative. They cancel over \(\F_2\), even before
normalization. Hence \(z*z=0\), proving the property.

Let \(c\in H_{m-1}(\Omega F)\) be the bottom class. Its image in
\(H_{m-1}K(m-1)\) has square zero for the additive multiplication.
This is the additive square-zero property just proved. It follows that
\(c^2\) lies in
the kernel of the projection to \(K(m-1)\). In degree \(2m-2\) this
kernel is one-dimensional, generated by the bottom class of the fibre.
The homotopy long exact sequence gives
\(\pi_{2m-2}\Omega F=\F_2\), and the space splitting shows that its
Hurewicz map identifies its nonzero element with that fibre class.

It remains to show that \(c^2\) is the nonzero fibre class. We do this
by showing that the product of bottom classes survives the cobar
spectral sequence. The space \(F\) is
\((m-1)\)-connected, so it is simply connected even when \(m=2\).
Adams's cobar construction \cite{Adams1956} has a multiplicative
comparison with chains on its based loop space; we use the modern
comparison in \cite[Theorem~1]{Rivera2019}. Choose a 1-reduced Kan
model for \(F\); the comparison for its cobar algebra follows from
\cite[Corollary~8 and Theorem~9]{Rivera2019}. Each desuspended
reduced chain has positive degree, so
the tensor-length filtration is finite in each total degree. The
resulting spectral sequence therefore converges to actual loop
homology, without a completion ambiguity. Write its first page as
\[
 E_1^{s,t}=(\widetilde H_*(F)^{\otimes s})_t,
 \qquad E_1^{s,t}\Longrightarrow H_{t-s}(\Omega F),
 \qquad d_r:(s,t)\longmapsto(s+r,t+r-1).
\]
Let \(c_0\in H_m(F)\) be the bottom class. In length two and
internal degree \(2m\), the term
\([c_0\mid c_0]\) is not a boundary for the coproduct differential:
the relevant coefficient is dual to the square of the bottom
cohomology class of \(F\), and that square vanishes by the defining
fibre construction. The term is a cycle since the reduced coproduct
of \(c_0\) is zero.

There is no higher incoming differential: its source would be
\((2-r,2m-r+1)\), hence have negative length or length zero and
positive internal degree. A higher outgoing differential
would land in length \(2+r\) and internal degree \(2m+r-1\), where
\(r\ge2\). Such a term vanishes, since each reduced homology factor
has degree at least \(m\) and
\[
(2+r)m>2m+r-1.
\]
This argument includes \(m=2\): the term is then in \((s,t)=(2,4)\),
and a higher outgoing target has internal degree \(r+3<2(r+2)\).
The same connectivity bound shows that \([c_0]\) is permanent and
detects the unique bottom class \(c\). Cobar concatenation represents
the loop product on the associated graded. The nonzero surviving
product \([c_0\mid c_0]\) consequently implies \(c^2\ne0\) in
actual loop homology. It is therefore the nonzero fibre class.

Under any lift \(T\to F\), the loop image of \(x\) is
\(\langle u,\sigma_*x\rangle c\), by homology suspension and the
one-dimensional bottom homology. Thus the adjoint of the composite
\(S^{2m-1}\to T\to F\) has Hurewicz image
\(\langle u,\sigma_*x\rangle c^2\). The injectivity just proved
shows that this composite is null precisely when the coefficient is zero.

Finally, \(\widetilde u^2=0\) is equivalent to the existence of a lift
of its representing map over \(C_g\) to \(F\). Such a lift restricts
to a lift on \(T\) whose composite with \(g\) is null; conversely,
any such lift extends over the attachment. The calculation applies to
every choice of lift. One can also see the independence directly:
changes of lift are measured by \(H^{2m-1}(T)\), whose pullback along
\(g\) is zero. This proves \eqref{eq:cone-square}.
\end{proof}

Retain the maps and classes of Lemma~\ref{lem:cone-square},
specialize to \(T=QS^n\) with \(n>0\), and suppose in addition
that \(x\) is primitive and \(A\)-annihilated.
Put \(B=\widetilde H^*(T;\F_2)\).
The homology Hopf algebra is polynomial on primitives, so its graded
dual cohomology algebra is exterior: each divided-power factor is
exterior on the generators at powers of two. Degreewise finite type
justifies the duality. In particular every positive cohomology class
has square zero. Define
\[
\lambda:B^m\longrightarrow\F_2,
\qquad \lambda(u)=\langle u,\sigma_*x\rangle.
\]
This functional vanishes on decomposables and on positive Steenrod
images, since \(\sigma_*x\) is primitive and \(A\)-annihilated.

\begin{proposition}\label{prop:cone-extension}
Under these assumptions, the cohomology extension of the cone induces
an extension of graded \(A\)-modules
\begin{equation}\label{eq:indecomposable-cone-extension}
0\longrightarrow\F_2\{a\}\longrightarrow\mathcal E
\longrightarrow Q B\longrightarrow0,
\end{equation}
where \(|a|=2m\). On any lift of a degree-\(m\) class \(v\in Q B\),
the operation \(\Sq^m\) has value \(\lambda(v)a\).
\end{proposition}

\begin{proof}
The hypothesis \(h(g)=0\) gives an exact sequence of positive
cohomology \(A\)-modules
\[
0\longrightarrow\F_2\{a\}\longrightarrow\widetilde H^*(C_g)
\longrightarrow B\longrightarrow0.
\]
The relative module structure implies that \(a\) times any
positive-degree class is zero. Choose exterior algebra generators of
\(B\) lifting a homogeneous basis of \(Q B\), choosing the basis in
degree \(m\) so that at most one generator has nonzero \(\lambda\).
Choose lifts of these generators to \(\widetilde H^*(C_g)\).
Their squares are zero except possibly for that one distinguished
degree-\(m\) generator, whose square is \(a\) by
Lemma~\ref{lem:cone-square}.

The nonempty square-free monomials in these lifts, together with \(a\), form a
graded vector-space basis of \(\widetilde H^*(C_g)\). Quotienting by
all products could kill \(a\), since it may itself be a square.
Instead we remove only products of distinct generators.
Let \(L\) be the span of the products of at least two distinct lifted
generators, with no repetitions. Restriction maps \(L\) isomorphically
onto \(B^2\). We claim that \(L\) is an \(A\)-submodule. Apply
Cartan to such a product and expand each factor in the lifted basis.
A term involving \(a\) vanishes against any remaining positive factor.
A repeated generator also gives zero, except that two occurrences of
the distinguished generator give \(a\); if another positive factor
remains, this term again vanishes.

The only possible contribution outside \(L\) is therefore a product
of exactly two linear occurrences of the distinguished generator.
A positive Steenrod operation on a factor cannot supply such a linear
coefficient, because \(\lambda\) annihilates positive Steenrod
images. If both operations are zero, the original two distinct
generators cannot both be the distinguished generator. This proves
the claim.

Taking the quotient by \(L\) gives
\eqref{eq:indecomposable-cone-extension}. In degree \(m\), the
square detector is unaffected, and its value is still \(\lambda\).
\end{proof}

\begin{corollary}\label{cor:tensor-detection}
If \(\lambda(v)=1\) for a degree-\(m\) monomial \(v\) in a
summand \(M_k(n)\) of \(Q B\), then
\begin{equation}\label{eq:tensor-detection}
\Sq^m\otimes v\ne0\quad\text{in }I\otimes_A M_k(n).
\end{equation}
The same conclusion holds for the bottom summand
\(M_0(n)=\Sigma^n\F_2\) when its degree is \(m\).
\end{corollary}

\begin{proof}
Pull back \eqref{eq:indecomposable-cone-extension} to the summand
\(M=M_k(n)\) and choose a graded additive section \(s\). In total
degree \(2m\), define the scalar \(\varphi(\alpha\otimes z)\) by
\[
\alpha s(z)-s(\alpha z)=\varphi(\alpha\otimes z)a,
\qquad \alpha\in I.
\]
It respects balancing. Indeed, in comparing
\(\alpha\beta\otimes z\) and \(\alpha\otimes\beta z\), the
intermediate degree \(|\beta|+|z|\) is strictly less than \(2m\)
because \(|\alpha|>0\). The possible discrepancy in the section
therefore vanishes, and associativity gives equality. Thus
\(\varphi\) descends to \((I\otimes_A M)_{2m}\).

Since \(n>0\), the module \(M_k(n)\) is a positive suspension of
an unstable polynomial module; hence \(\Sq^m v=0\). This is also
true in the one-dimensional bottom summand. Proposition~\ref{prop:cone-extension}
therefore gives \(\varphi(\Sq^m\otimes v)=\lambda(v)=1\), proving
\eqref{eq:tensor-detection}.
\end{proof}
